% Part: normal-modal-logic
% Chapter: syntax-and-semantics
% Section: language-modal-logic

\documentclass[../../../include/open-logic-section]{subfiles}

\begin{document}

\olfileid{nml}{syn}{lan}

\olsection{基本模态逻辑的语言}

\begin{defn}
基本模态逻辑的语言包含：
\begin{enumerate}
  \tagitem{prvFalse}{表示假的命题常项~$\lfalse$。}{}
  \tagitem{prvTrue}{表示真的命题常项~$\ltrue$。}{}
  \item 一个可数无穷的命题变元集合：$\Obj p_0$, $\Obj p_1$, $\Obj p_2$, \dots
  \item 命题联结词： \startycommalist
  \iftag{prvNot}{\ycomma $\lnot$ (否定)}{}%
  \iftag{prvAnd}{\ycomma $\land$ (合取)}{}%
  \iftag{prvOr}{\ycomma $\lor$ (析取)}{}%
  \iftag{prvIf}{\ycomma $\lif$ (条件句)}{}%
  \iftag{prvIff}{\ycomma $\liff$ (双条件句)}{}。
  \tagitem{prvBox}{模态算子 $\Box$。}{}
  \tagitem{prvDiamond}{模态算子 $\Diamond$。}{}
\end{enumerate}
\end{defn}

\begin{defn}
基本模态语言的\emph{公式}归纳定义如下：
\begin{enumerate}
\tagitem{prvFalse}{$\lfalse$ 是一个原子公式。}{}

\tagitem{prvTrue}{$\ltrue$ 是一个原子公式。}{}

\item 每个命题变元 $\Obj p_i$ 是一个（原子）公式。

\tagitem{prvNot}{若 $!A$ 是公式，则 $\lnot !A$ 是公式。}{}

\tagitem{prvAnd}{若 $!A$ 和 $!B$ 是公式，则 $(!A \land !B)$ 是公式。}{}

\tagitem{prvOr}{若 $!A$ 和 $!B$ 是公式，则 $(!A \lor !B)$ 是公式。}{}

\tagitem{prvIf}{若 $!A$ 和 $!B$ 是公式，则 $(!A \lif !B)$ 是公式。}{}

\tagitem{prvIff}{若 $!A$ 和 $!B$ 是公式，则 $(!A \liff !B)$ 是公式。}{}

\tagitem{prvBox}{若 $!A$ 是公式，则 $\Box !A$ 是公式。}{}

\tagitem{prvDiamond}{若 $!A$ 是公式，则 $\Diamond !A$ 是公式。}{}

\tagitem{limitClause}{除此之外皆非公式。}{}
\end{enumerate}
\end{defn}

\begin{tagblock}{defNot,defOr,defAnd,defIf,defIff,defTrue,defFalse,defBox,defDiamond}
\begin{defn}
使用被定义的算子构造的公式应按如下方式理解：

\begin{tagenumerate}{defTrue,defFalse,defNot,defOr,defAnd,defIf,defIff,defBox,defDiamond}

\tagitem{defTrue}{$\ltrue$ 是
  \iftag{prvFalse}{$\lnot\lfalse$}{$(!A \lor \lnot !A)$，其中 $!A$ 为某个固定的原子公式}的缩写。}{}

\tagitem{defFalse}{$\lfalse$ 是
  \iftag{prvTrue}{$\lnot\ltrue$}{$(!A \land \lnot !A)$，其中 $!A$ 为某个固定的原子公式}的缩写。}{}

\tagitem{defNot}{$\lnot !A$ 是 $!A \lif \lfalse$ 的缩写。}{}

\tagitem{defOr}{$!A \lor !B$ 是
  \iftag{prvAnd}{$\lnot(\lnot !A \land \lnot !B)$}{$\lnot !A \lif !B$}的缩写。}{}

\tagitem{defAnd}{$!A \land !B$ 是
  \iftag{prvOr}{$\lnot(\lnot !A \lor \lnot !B)$}{$\lnot (!A \lif \lnot !B)$}的缩写。}{}

\tagitem{defIf}{$!A \lif !B$ 是
  \iftag{prvOr}{$\lnot !A \lor !B)$}{$\lnot (!A \land \lnot !B)$}的缩写。}{}

\tagitem{defIff}{$!A \liff !B$ 是 $(!A \lif !B) \land (!B
  \lif !A)$ 的缩写。}{}

\tagitem{defBox}{$\Box !A$ 是 $\lnot\Diamond\lnot !A$ 的缩写。}{}

\tagitem{defDiamond}{$\Diamond !A$ 是 $\lnot\Box\lnot !A$ 的缩写。}{}
\end{tagenumerate}
\end{defn}
\end{tagblock}

若公式~$!A$ 不包含 $\Box$ 或 $\Diamond$，则称它是\emph{模态自由的}。

\end{document}